\section{Κίνητρο και Σκοπός της Εργασίας}
\label{sec:intro_motivation}

Ο αεροδυναμικός σχεδιασμός ενός αεροσκάφους ---και ειδικότερα ενός 3D-εκτυπώσιμου μη επανδρωμένου οχήματος (UAV)--- αποτελεί ένα σύνθετο πολυπαραμετρικό πρόβλημα βελτιστοποίησης~\cite{goldberg1989genetic}. Η παραδοσιακή αξιολόγηση μέσω Υπολογιστικής Ρευστοδυναμικής (Computational Fluid Dynamics - CFD) είναι υπολογιστικά εξαιρετικά βαριά. Μία μεμονωμένη προσομοίωση Navier-Stokes μπορεί να απαιτήσει από αρκετά λεπτά έως ολόκληρες ώρες υπολογιστικού χρόνου.

Όταν η μέθοδος βελτιστοποίησης βασίζεται σε εξελικτικούς αλγορίθμους (Genetic Algorithms - GAs), απαιτούνται χιλιάδες ή δεκάδες χιλιάδες διαδοχικές αξιολογήσεις υποψηφίων σχεδίων~\cite{jin2011surrogate}. Η εκτέλεση ενός κλασικού γενετικού αλγορίθμου υπό αυτές τις συνθήκες καθίσταται απαγορευτική σε πρακτικά χρονικά όρια. Συνεπώς, γεννάται η επιτακτική ανάγκη για έξυπνες αρχιτεκτονικές που μειώνουν δραστικά τις κλήσεις προς τον βαρύ προσομοιωτή, αξιοποιώντας τεχνικές κατανεμημένης αποθήκευσης και προσεγγιστικής μάθησης (surrogate modeling).
